\documentclass[10pt]{article}
\providecommand{\FIGDIR}{fig}
\newcommand{\DIGESTDATE}{30 October 2026}
\input{preamble}

\begin{document}

% ==================================================================== MASTHEAD
\thispagestyle{empty}
\begin{tikzpicture}[remember picture, overlay]
  \fill[Deep] (current page.north west) rectangle
      ($(current page.north east)+(0,-67mm)$);
  \fill[Amber] ($(current page.north west)+(0,-67mm)$) rectangle
      ($(current page.north east)+(0,-68.4mm)$);
  \foreach \x/\y/\r in {22/12/0.9, 41/26/1.5, 63/9/0.7, 88/31/1.1, 112/17/1.9,
                        138/29/0.8, 159/13/1.3, 178/54/1.0, 196/21/1.6,
                        30/44/0.6, 74/57/1.2, 125/48/0.9, 168/6/0.8, 191/59/0.7,
                        52/18/1.0, 100/60/0.6, 148/20/1.1, 66/49/1.4, 205/45/0.8}
    \fill[Sky, opacity=0.30] ($(current page.north west)+(\x mm,-\y mm)$) circle (\r mm);
\end{tikzpicture}

\vspace*{4mm}
{\sffamily\fontsize{7.6}{9}\selectfont\color{Sky}%
 AUTOMATED DAILY LITERATURE DIGEST\;\textbullet\;PhD RESEARCH WATCH\par}
\vspace{2.2mm}
{\sffamily\bfseries\fontsize{27}{29}\selectfont\color{white}%
 Particulate Matter\par}
{\sffamily\fontsize{27}{29}\selectfont\color{Sky}%
 Monitoring \& Health Effects\par}

\vspace{5.0mm}
{\sffamily\fontsize{8.4}{10}\selectfont\color{white}%
 Issue 2026-10-30\;\;\textcolor{Amber}{\textbar}\;\;Entry window 30 October 2026 (full day)\;\;%
 \textcolor{Amber}{\textbar}\;\;24 records screened in scope\par}
\vspace{18mm}

\noindent
\begin{minipage}[t]{0.190\linewidth}\metric{24}{RECORDS IN SCOPE\\(67 EXCLUDED)}{Deep}\end{minipage}\hfill
\begin{minipage}[t]{0.190\linewidth}\metric{9}{OF 10 SUBTOPIC\\BANDS POPULATED}{Teal}\end{minipage}\hfill
\begin{minipage}[t]{0.190\linewidth}\metric{3}{EFFECT ESTIMATES\\WITH A CI}{Sky}\end{minipage}\hfill
\begin{minipage}[t]{0.190\linewidth}\metric{1}{TIER-A\\RECORDS}{Amber}\end{minipage}\hfill
\begin{minipage}[t]{0.190\linewidth}\metric{5}{RECORDS CARRIED ON\\METADATA ALONE}{Clay}\end{minipage}

\vspace{6mm}

\section*{\sffamily\bfseries\color{Deep}Signal of the day}
\vspace{-1.5mm}{\color{Amber}\rule{\linewidth}{1.1pt}}\vspace{2.5mm}

\begin{keybox}
{\sffamily\bfseries\small\color{Deep}All of Us, 179,218 adults, preregistered: heat anomaly and social
vulnerability predict incident depression; 30-day wildfire-smoke PM$_{2.5}$ does not --- at 3-digit-ZIP resolution.}\\[1.4mm]
\footnotesize
Cox models on 5,484 first EHR-recorded depression diagnoses (mean follow-up 371 days). Heat anomaly HR 1.06 (1.03--1.10)
and 1.08 (1.04--1.11) per SD; social-material vulnerability HR 1.36 (1.33--1.40) per SD. Mean smoke PM$_{2.5}$ (Stanford
ECHO): no association in the primary model; secondary smoke metrics heterogeneous by metric and follow-up window.
\textbf{Where it is weak:} exposure is a ZIP3 average frozen at baseline --- smoke episodes are short and spatially
patchy, so this is the textbook setting for attenuation to the null. A preregistered null is still a null on a
heavily smoothed exposure, not evidence of no effect.
\textbf{Why it matters here:} the contrast between a clean heat signal and a flat smoke signal is what residence-scale,
humidity-corrected sensor PM$_{2.5}$ would test directly.\par
\end{keybox}

\vspace{3mm}

\section*{\sffamily\bfseries\color{Deep}What changed today}
\vspace{-1.5mm}{\color{Amber}\rule{\linewidth}{1.1pt}}\vspace{2.5mm}

\footnotesize
\begin{itemize}

\item \textbf{Low-cost sensors as the reference.} A MethodsX protocol measures air-cleaner CADR from PM$_{2.5}$ decay on
low-cost sensors in a real home: $\sim$26\,\% below lab values from research-grade instruments --- conservative, and
quantified. A four-filter Corsi--Rosenthal box led the DIY builds.

\item \textbf{Hygroscopicity is source-dependent.} Eastern China, humidified nephelometry: $\kappa_{\rm OA}$ 0--0.49
(mean 0.11), tracking biomass-burning influence rather than O:C. A fixed-$\kappa$ humidity correction is wrong during
smoke. Separately, Kelvin-suppressed growth of 15--36\,nm particles changes modelled lung deposition by up to 25\,\%.

\item \textbf{Policy co-benefit, modelled.} China's 2017 climate-adaptive city pilot: PM$_{2.5}$ $-$3.87\,$\mu$g\,m$^{-3}$
($-$6.49 to $-$1.25) by difference-in-differences, $\sim$8\,\% lower GEMM burden --- shrinking in the matched sample.

\item \textbf{Prenatal exposure, small signals.} TTN in Taiwan: PM$_{2.5}$ aOR 1.001 (1.000--1.002) per unit, dominated by
CO; ECHO cord-blood methylation classifies high PM$_{2.5}$ exposure with AUC only 0.64--0.68.

\item \textbf{Five metadata-only records (21\,\%)}, two of them squarely in scope for sensing: PN-PTI counters on gasoline
exhaust (\emph{J.\ Aerosol Sci.}) and e-cigarette aerosol decay. Retry first.

\end{itemize}

\vspace{2mm}

\begin{keybox}
{\sffamily\bfseries\footnotesize\color{Deep}Window continuity}\\[1.2mm]
{\fontsize{7.2}{8.8}\selectfont
The 29 October issue closed with \texttt{last\_entry\_date} at 29 October. A 31 October 05:33 UTC firing harvested
the PubMed and Europe PMC legs for 30 October and stopped. This issue was built on the \textbf{1 November 05:25 UTC
(00:25 CDT) run}; by the 22:00 rule the newest buildable day is 31 October, so the run backfills \textbf{30 October in
full} here and 31 October as its own issue, with the W44 weekly (25--31 October) and the October monthly. Candidates
were screened sequentially across the two days (cross-day DOI and title set built before authoring) to avoid the
double admissions of the 23--29 October backfill.\par}
\end{keybox}

\vspace{3mm}
\par\vskip 0pt plus 18\baselineskip\penalty-200\vskip 0pt plus -18\baselineskip
\section*{\sffamily\bfseries\color{Deep}Corpus analytics}
\vspace{-1.5mm}{\color{Amber}\rule{\linewidth}{1.1pt}}\vspace{2.5mm}

\noindent\begin{minipage}{\linewidth}\includegraphics[width=\linewidth]{\FIGDIR/f1_subtopics.png}
\figcap{Subtopic assignment of the 24 in-scope records across nine bands; other clinical endpoints is empty.
Sensing and exposure assessment lead with five each, mechanistic toxicology four; reproductive, respiratory, burden and
occupational/indoor two each; neuro and cardiovascular one each.}\end{minipage}

\vspace{2mm}
\noindent\begin{minipage}{\linewidth}\includegraphics[width=\linewidth]{\FIGDIR/f2_design_endpoint.png}
\figcap{Left --- architecture: \textbf{metadata-only 5 (21\,\%)}, four measurement campaigns, four
toxicology, three modelling, three cross-sectional, two cohort-type (All of Us, TTN case-cohort), two ecological and one
laboratory (HTDMA). Right --- endpoints span depression, TTN, cord-blood methylation, air-cleaner CADR,
hygroscopicity and new particle formation.}\end{minipage}

\vspace{2mm}
\noindent\begin{minipage}{\linewidth}\includegraphics[width=\linewidth]{\FIGDIR/f3_metric_geography.png}
\figcap{Left --- particle metric: PM$_{2.5}$-only 10, unspeciated/emissions 6, size-resolved/number
3, optical 3, and one each joint PM$_{2.5}$+PM$_{10}$ and composition. Right --- geography: global/not stated 5, North
America 4, China 4, Europe 4, South Asia 3, and one each East Asia (Taiwan), Sub-Saharan Africa (Tanzania), Latin
America (Brazil) and Central Asia (Kazakhstan).}\end{minipage}

\vspace{2mm}
\noindent\begin{minipage}{\linewidth}\includegraphics[width=\linewidth]{\FIGDIR/f6_lifecourse.png}
\figcap{Life-course windows: three in-utero/early-life records (TTN, ECHO cord blood, CHAMACOS prenatal
exposure), one childhood (Dar es Salaam school BP), one adolescent/young-adult (nasal EWAS), one working-age (All of Us
depression); no older-adult record today.}\end{minipage}

\vspace{2mm}
\noindent\begin{minipage}{\linewidth}\includegraphics[width=\linewidth]{\FIGDIR/f4_heatmap.png}
\figcap{Subtopic $\times$ study architecture. Sensing splits between measurement, modelling and metadata
cells; mechanistic toxicology is all experimental; the health bands sit in cohort, cross-sectional and ecological
cells.}\end{minipage}

\vspace{2mm}

\noindent\begin{minipage}{\linewidth}\includegraphics[width=\linewidth]{\FIGDIR/f5_forest.png}
\figcap{Ratio scale, one estimate: prenatal PM$_{2.5}$ and transient tachypnea of the newborn (aOR per unit
concentration, Taiwan, 3,091 cases). The CI touches 1; CO, not PM$_{2.5}$, carries the study.}\end{minipage}

\vspace{3mm}

\noindent\begin{minipage}{\linewidth}\includegraphics[width=\linewidth]{\FIGDIR/f5_forest_2.png}
\figcap{Difference scale, one study (China, 264 cities): the climate-adaptive city pilot's PM$_{2.5}$ change
($\mu$g\,m$^{-3}$) and the GEMM-modelled burden change (\%). Different units; the second is derived from the first and
is not an independent estimate.}\end{minipage}

\vspace{4mm}

\par\vskip 0pt plus 24\baselineskip\penalty-200\vskip 0pt plus -24\baselineskip
\section*{\sffamily\bfseries\color{Deep}Paper-level digest}
\vspace{-1.5mm}{\color{Amber}\rule{\linewidth}{1.1pt}}\vspace{2.5mm}

{\fontsize{7.4}{9}\selectfont\color{Slate}Ordered by cluster. Each entry gives the
methodological core and the finding that matters, not an abstract paraphrase. Records
marked \textit{metadata only} carry no recoverable abstract and assert no finding.\par}

\band{Neuro / mental health\hfill 1 record}{Deep}

\paperentry{Yang heat anomaly wildfire smoke PM$_{2.5}$ incident depression All of Us 2026}
{Heat anomaly, wildfire smoke, social-material vulnerability, and incident depression in a national EHR-linked cohort: A preregistered secondary analysis}
{J Affect Disord}
{Preregistered Cox analysis in All of Us: 179,218 adults free of recorded depression, 5,484 incident cases over a mean 371 days. Exposures were the 30-day pre-index PRISM maximum-temperature anomaly, Stanford ECHO smoke PM$_{2.5}$ and a social-material vulnerability (SMV) composite. Heat anomaly: HR 1.06 (1.03-1.10) and 1.08 (1.04-1.11) per SD; SMV: HR 1.36 (1.33-1.40) per SD. Mean smoke PM$_{2.5}$ showed no association in the primary model; secondary smoke metrics were heterogeneous by metric and follow-up window. Limitation: exposure is assigned at the 3-digit ZIP level and frozen at baseline, so smoke-episode contrast is heavily smoothed and a null is weak evidence of no effect; EHR ascertainment depends on care access. Relevance: high - this is the misclassification problem that residence-scale sensor-corrected smoke PM$_{2.5}$ exists to fix.}
{10.1016/j.jad.2026.122663}

\par\vskip 0pt plus 10\baselineskip\penalty-200\vskip 0pt plus -10\baselineskip
\band{Cardiovascular \& metabolic\hfill 1 record}{Teal}

\paperentry{Musa outdoor air pollution cardiovascular health school children Dar es Salaam 2026}
{Effects of outdoor air pollution on the cardiovascular health of school-aged children residing near polluted areas and the efficacy of mitigation measures in Dar es Salaam, Tanzania}
{BMC Public Health}
{Community cross-sectional study, December 2024-July 2025, of 653 school-aged children living in areas classed as highly polluted by proximity to industrial zones, major roads and waste-burning sites. Single oscillometric BP reading; 19.9\% had elevated BP. In log-binomial models, prior respiratory disease carried ARR 8.81 (6.12-12.67) for cardiovascular symptoms. Awareness of mitigation policy was mixed. Limitation: no pollutant was measured and there is no unexposed comparison group, so the title's causal framing is not supported; one-visit BP overstates prevalence in children. Relevance: high as a gap marker - urban East Africa has almost no PM monitoring, and this design shows what health studies look like without it.}
{10.1186/s12889-026-28148-9}

\par\vskip 0pt plus 10\baselineskip\penalty-200\vskip 0pt plus -10\baselineskip
\band{Respiratory \& allergic\hfill 2 records}{Sky}

\paperentry{Perez-Garcia nasal EWAS respiratory biomarkers Mexican American CHAMACOS 2026}
{Nasal Epigenome-Wide Association Study of Respiratory and Allergic Biomarkers in Mexican American Young Adults}
{medRxiv}
{Preprint. Nasal EWAS (EPICv2, 804,220 CpGs) in 354 CHAMACOS young adults aged 22, replicated in 540 Project Viva adolescents aged 13. Forty-nine CpGs were associated with FeNO and 557 DMRs with respiratory outcomes (FDR<0.05); 26 CpGs and 103 DMRs replicated (e.g. LRRFIP1, EPX, IL13). Ten FeNO-related CpGs were associated with prenatal and/or past-year PM$_{2.5}$ or O$_3$ at nominal p<0.05 only. Limitation: the air-pollution step is a secondary look-up without multiple-testing control, so the PM$_{2.5}$ link is hypothesis-generating; cross-sectional nasal sampling cannot order exposure and methylation. Relevance: low-medium - an agricultural-valley cohort where PM sources are local and poorly served by regulatory monitors.}
{10.64898/2026.10.27.26366640}

\paperentry{Kobza respiratory ambulatory care air quality Upper Silesia 2026}
{The assessment of respiratory diseases in ambulatory care attributable to air quality in the Upper Silesian agglomeration}
{Int J Occup Med Environ Health}
{Five-year ecological analysis correlating daily mean concentrations from about 50 monitoring stations with National Health Fund outpatient incidence for seven respiratory ICD codes across eight regions of the Upper Silesian voivodeship. The strongest signal was for SO$_2$, not PM: tau 0.99 and rho 0.9 in regions 7 and 8, rho 0.78 in regions 1 and 4. Limitation: rank correlations on aggregated series without control for season, influenza, temperature or autocorrelation - seasonality alone would generate these coefficients, so the results carry no causal or PM-specific information. Relevance: low - included because PM$_{2.5}$/PM$_{10}$ were analysed; the coal-heating region is a classic case where denser residential PM sensing would sharpen exposure contrast.}
{10.13075/ijomeh.1896.02875}

\par\vskip 0pt plus 10\baselineskip\penalty-200\vskip 0pt plus -10\baselineskip
\band{Reproductive \& developmental\hfill 2 records}{Sage}

\paperentry{Lin prenatal air pollution transient tachypnea newborn Taiwan 2026}
{Prenatal air pollution exposure and the risk of transient tachypnea of the newborn}
{Front Pediatr}
{Nationwide matched design on the Taiwan NHIRD linked to the Maternal and Child Health Database, 2004-2020: 3,091 TTN cases in first-born infants and 12,364 controls (1:4 on maternal age, sex, birth date). Exposure over the 40 weeks before birth from the nearest of 71 stations to the residential postal code. Adjusted ORs: CO 1.806 (1.312-2.487), NOx 1.010 (1.005-1.015), PM$_{2.5}$ 1.001 (1.000-1.002), the last per unit concentration with a CI touching the null. Effects concentrated in the Taipei basin. Limitation: nearest-station assignment at postal-code scale, no multipollutant model reported, and the PM$_{2.5}$ signal is marginal while CO dominates, pointing to traffic as a mixture rather than PM. Relevance: medium - traffic-gradient exposure in a dense basin is where street-level sensor networks change exposure contrast.}
{10.3389/fped.2026.1941575}

\paperentry{Gaylord DNA methylation signature prenatal air pollution ECHO cohort 2026}
{A methodology to identify a DNA methylation-based signature of prenatal air pollution exposure in the ECHO Cohort}
{Environ Epigenet}
{Refines a pipeline that predicts high versus low in-utero air-pollution exposure from cord-blood DNA methylation (450K/EPIC) across multiple ECHO pregnancy cohorts. Exposures were ML-modelled at geocoded addresses; training and validation sets were matched on exposure level, site and array. Validation AUCs: PM$_{2.5}$ 0.64-0.68, NO$_2$ 0.62-0.95, O$_3$ 0.50-0.75, varying with region, exposure window and the high/low threshold. Limitation: the target is itself a model output, so the epigenetic classifier inherits exposure-model error, and PM$_{2.5}$ discrimination is modest; dichotomisation discards dose information. Relevance: medium - a biomarker of cumulative exposure is only as good as the exposure surface it is trained on, which is where denser monitoring would raise the ceiling.}
{10.1093/eep/dvag034}

\par\vskip 0pt plus 10\baselineskip\penalty-200\vskip 0pt plus -10\baselineskip
\band{Mechanistic toxicology\hfill 4 records}{Violet}

\paperentry{He FOXM1 lncRNA COPD-R biomass PM$_{2.5}$ pulmonary inflammation rat 2026}
{FOXM1/LncRNA COPD-R regulates pulmonary inflammation in chronic obstructive pulmonary disease induced by biomass-related PM$_{2.5}$}
{Ecotoxicol Environ Saf}
{Rats exposed to biomass-fuel smoke for seven months developed COPD-like disease; lung RNA-seq found 312 differentially expressed lncRNAs (294 up). lncRNA COPD-R rose from 0.15 to 1.05 (relative expression, p<0.01). In rat L2 alveolar cells, biomass-related PM$_{2.5}$ increased COPD-R; knockdown reduced PM-induced CCL2, IL-17F, MMP-9 and TIMP-1, and ChIP/luciferase assays placed FOXM1 upstream as its transcriptional activator, with rescue experiments consistent with a FOXM1-COPD-R axis. Limitation: whole-smoke exposure in vivo but PM$_{2.5}$ only in vitro, dose not reported in the abstract, and a rat lncRNA whose human orthologue is not established. Relevance: low - mechanistic support for household biomass PM, where the exposure gap is measurement in homes.}
{10.1016/j.ecoenv.2026.120952}

\paperentry{Chen pirfenidone fine particulate matter EMT mitochondrial fission autophagy 2026}
{Pirfenidone reduces fine particulate matter-induced epithelial to mesenchymal transition in lung epithelial cells through regulating mitochondrial fission and autophagy}
{Respir Res}
{A549 cells exposed to 100 ug/mL fine PM lost E-cadherin and gained vimentin, alpha-SMA, TGF-beta1, Snail and Slug, turned spindle-shaped and migrated more; PM increased mitochondrial fission and LC3B/p62 accumulation, i.e. impaired autophagic flux. Intratracheal PM in mice reproduced EMT markers, inflammatory infiltration and alveolar wall thickening. Pirfenidone reversed these effects in vitro and reduced collagen deposition in vivo. Limitation: 100 ug/mL on a cell monolayer and bolus instillation are far above realistic deposited doses, and A549 is a carcinoma line; the drug result says nothing about population exposure. Relevance: low - mechanism for PM-associated fibrosis; no measurement content.}
{10.1186/s12931-026-03804-1}

\paperentry{Zakari PM$_{1}$ microbial profiles lung gut microbiome mice 2026}
{Associations of particulate matter-associated microbial profiles with lung and intestinal microbiomes in healthy mice}
{J Hazard Mater}
{Mice were housed under ambient PM$_{1}$, HEPA-filtered air or clean air; PM$_{1}$ was characterised for ions, metals, carbon fractions, oxidative potential and its own 16S microbiome (Proteobacteria-dominated). Lung diversity barely changed while the gut microbiome shifted: PM$_{1}$ reduced Actinobacteriota and Bacteroidota and raised Proteobacteria versus both control arms. Eleven families and six genera overlapped among PM$_{1}$, lung and gut communities in the PM$_{1}$ group. Limitation: group sizes and exposure concentrations are not given in the abstract, gases co-vary with PM$_{1}$, and taxon overlap does not show transfer from air to gut. Relevance: low - airborne microbial load is a PM component no optical sensor resolves.}
{10.1016/j.jhazmat.2026.143978}

\paperentry{Gautam diesel exhaust miR-133b-3p CTGF Hippo ECM remodeling mouse lung 2026}
{Diesel Exhaust-Induced miR-133b-3p Downregulation and CTGF-Linked Hippo Pathway Inactivation Promote ECM Remodeling and Inflammation in Mouse Lungs}
{Environ Mol Mutagen}
{C57BL/6 mice received diesel exhaust 30 min/day, 5 days/week for 8 weeks versus filtered air. Microarrays showed miR-133b-3p down more than 2-fold; CTGF was confirmed as a direct target in A549 cells (mimic reduced CTGF mRNA p=0.0054, protein p=0.0001). In vivo, exhaust raised CTGF, lowered P-YAP/YAP (Hippo inactivation), increased BALF IL-18, IL-1beta and GM-CSF, raised MMP-3/7/12 and lowered LOXL1/4, consistent with ECM remodelling. Limitation: exhaust dose (PM mass, NO$_2$) is not reported in the abstract, so particle and gas effects cannot be separated; the CTGF link is correlative in vivo. Relevance: low - traffic-exhaust mechanism; supports particle number or black carbon over mass as the near-road metric.}
{10.1002/em.70088}

\par\vskip 0pt plus 10\baselineskip\penalty-200\vskip 0pt plus -10\baselineskip
\band{Exposure assessment \& modelling\hfill 5 records}{Clay}

\paperentry{Yousaf nanoscale curvature hygroscopic growth respiratory deposition ultrafine 2026}
{Nanoscale Curvature Controls Hygroscopic Growth and Respiratory Deposition of Atmospheric Nanoparticles}
{Environ Pollut}
{HTDMA measurements of 15-100 nm NaCl particles fed into the MPPD lung-deposition model. The Kelvin effect raised deliquescence RH to 77-81\% for 15-36 nm particles and suppressed growth; at 80-90\% RH, growth factors were 1.95-2.25 (15 nm) and 1.99-2.36 (100 nm). Ignoring hygroscopic growth overestimated total deposition by up to 25\%, with an 18\% tracheobronchial difference for particles <=50 nm, and shifted larger UFPs deeper into the lung. Limitation: pure NaCl is a best-case hygroscopic aerosol; ambient UFP is organic-rich with far lower kappa, so the correction is an upper bound. Relevance: medium - the same size-dependent water uptake that biases optical sensors in humid air also changes delivered dose, so dry-size number counts understate the problem.}
{10.1016/j.envpol.2026.129429}

\paperentry{Kalkavouras ammonia sulfuric acid BVOC urban new particle formation Athens 2026}
{Assessing the combined roles of ammonia, sulfuric acid, and biogenic volatile organic compounds in urban new particle formation}
{Environ Pollut}
{Seven months (June 2023-January 2024) of particle number size distributions, NH$_3$ and BVOCs at the Thissio urban background site in Athens, with a proxy for H2SO4. Mean NH$_3$ 2.3 +/- 0.7 $\mu$g\,m$^{-3}$. NPF occurred on 19\% of days, most often in autumn (29-33\%) when ventilation and condensation sink were low. Precursor levels barely differed between event and non-event days; the highest J14 came at intermediate NH$_3$ and H2SO4 with low BVOCs, while growth rates rose with NH$_3$ and isoprene/monoterpenes. Limitation: H2SO4 is a proxy, not measured, and the record lacks a full annual cycle. Relevance: medium - NPF bursts dominate urban number concentration but are invisible to mass-based optical sensors, which see nothing below \textasciitilde{}300 nm.}
{10.1016/j.envpol.2026.129453}

\paperentry{Chakraborty fog intensity droplet size aqueous SOA Kanpur 2026}
{Fog Intensity and Droplet Size-Resolved Aqueous Secondary Organic Aerosol (aqSOA) Formation in Polluted Winter Fog: New Insights from a Combined Observational and Modeling Study}
{Environ Sci Technol}
{Two winters of size-resolved fogwater at Kanpur analysed by offline HR-AMS and PMF; three highly oxidised factors (O/C 0.95-1.19) were summed as an aqSOA proxy and scaled by liquid water content and WSOC. Across seven events aqSOA was 3.1-8.4 $\mu$g\,m$^{-3}$, a median \textasciitilde{}7\% of ambient SOA and 5\% of OA, up to 13\% of SOA in the densest fogs. Fine droplets carried 43-84\% of aqSOA, linked to dissolved Fe and Cu and longer residence. Limitation: the aqSOA proxy is a PMF construct, and offline AMS of fogwater may not represent residual particle mass after evaporation. Relevance: medium - Indo-Gangetic fog episodes are when optical PM sensors fail outright from droplet interference, while secondary mass is being added.}
{10.1021/acs.est.6c09700}

\paperentry{Yingkar PM$_{2.5}$-bound PAHs Central Asia urban trends sources risk 2026}
{Environment-to-Health Linkages of PM$_{2.5}$-Bound Polycyclic Aromatic Hydrocarbons in Urban Cities of Central Asia: Trends, Sources, and Risk Assessment}
{Atmos Environ}
{\textsc{metadata only.} Title and authors only (Yingkar, Omarova, Sovetova, Alibekov et al.). No abstract was retrievable from Crossref or Semantic Scholar at harvest. The title describes trends, source apportionment and health-risk assessment of PM$_{2.5}$-bound PAHs across urban Central Asian cities - a region with sparse speciation data and heavy coal heating. Cities, sampling design, concentrations and risk estimates are not available; nothing is inferred. Relevance: medium - coal-heating PAH peaks in Almaty-type basins are where low-cost PM networks already operate without any chemical reference; retry the abstract on a later run.}
{10.1016/j.atmosenv.2026.122447}

\paperentry{Fang end of uniformity spatial climatic toxicological redefinitions air quality 2026}
{The End of Uniformity: Spatial, Climatic, and Toxicological Redefinitions of Air Quality}
{Environ Sci Technol}
{\textsc{metadata only.} Title and authors only (Fang, Song, Weon). No abstract was retrievable at harvest; the format is likely a viewpoint or perspective. The title argues against uniform air-quality definitions along three axes - spatial heterogeneity, climate-driven change and toxicity beyond mass - which maps directly onto hyperlocal sensing, wildfire-era exposure and component- or oxidative-potential-based metrics. Arguments and any proposals are not available; nothing is inferred. Relevance: medium - if it advocates spatially resolved or toxicity-weighted standards, it is a policy argument for distributed networks; retry the text on a later run.}
{10.1021/acs.est.6c14170}

\par\vskip 0pt plus 10\baselineskip\penalty-200\vskip 0pt plus -10\baselineskip
\band{Sensing, forecasting \& instrumentation\hfill 5 records}{Amber}

\paperentry{Waghmare machine learning explainable AI PM$_{2.5}$ prediction Nagpur 2026}
{Machine learning models with explainable AI for PM$_{2.5}$ prediction for Urban air quality assessment}
{MethodsX}
{Six regressors (decision tree, extra trees, random forest, KNN, XGBoost, SVR) trained on \textasciitilde{}4,000 observations of PM$_{2.5}$, NO$_2$, SO$_2$, NH$_3$, O$_3$ and meteorology in Nagpur, 70/30 split with 10-fold CV. A single decision tree scored best: test R2 0.90, RMSE 7.5-8.9 $\mu$g\,m$^{-3}$, MAE 5.4-6.6 $\mu$g\,m$^{-3}$; SHAP ranked NO$_2$, O$_3$, RH and NH$_3$ top. Limitation: random splits of autocorrelated time series leak information, so skill is inflated; co-pollutants measured at the same station make this a nowcast, not a forecast, and a lone tree beating ensembles is a warning sign. Relevance: low-medium - illustrates the validation discipline (blocked or spatial CV) that sensor-network ML calibration needs.}
{10.1016/j.mex.2026.104217}

\paperentry{Wu source-dependent organic aerosol hygroscopicity eastern China biomass burning 2026}
{Source-dependent organic aerosol hygroscopicity in eastern China: observational evidence for conditional enhancement under biomass-burning influence}
{Atmos Chem Phys}
{Winter field campaign at a rural eastern-China site retrieving organic-aerosol hygroscopicity from humidified nephelometry and ZSR closure. kappa\_OA ranged from \textasciitilde{}0 to 0.49 (mean 0.11 +/- 0.11) and correlated only weakly with O:C; instead it rose with biomass-burning influence, and ageing enhanced kappa\_OA only when biomass burning dominated. Highly oxidised OA under weak biomass influence stayed weakly hygroscopic. Limitation: kappa\_OA is a residual after subtracting inorganic contributions, so errors in the inorganic closure load onto the organic term; one site, one season. Relevance: high - humidity correction of optical PM sensors usually assumes a fixed kappa; this shows kappa shifts with source, so smoke-affected periods need a different growth correction.}
{10.5194/acp-26-15333-2026}

\paperentry{Schmidt AOD gap filling iterative singular spectrum analysis AERONET Cuiaba 2026}
{Gap filling in long-term AOD time series using Iterative Singular Spectrum Analysis at a tropical AERONET site}
{Environ Monit Assess}
{Applies iterative singular spectrum analysis (ISSA) to fill a 2005-2024 daily AOD record from the Cuiaba-Miranda AERONET station, where \textasciitilde{}47\% of days are missing, mostly in the cloudy wet season. Climatological (day-of-year mean) initialisation beat global-mean initialisation, giving lower error and better temporal consistency, and the cross-validation design materially changed the optimal reconstruction parameters. No error magnitudes are given in the abstract. Limitation: gaps are not missing at random - they coincide with cloud and wet-season aerosol regimes - so a climatology-anchored fill will bias toward dry-season behaviour exactly where data are absent. Relevance: low-medium - the same informative-missingness problem affects sensor networks with humidity-triggered dropouts and satellite-surface fusion.}
{10.1007/s10661-026-16067-6}

\paperentry{Trivanovic periodic technical inspection particle number devices gasoline exhaust 2026}
{Assessment of existing PN-PTI devices to measure gasoline exhaust}
{J Aerosol Sci}
{\textsc{metadata only.} Title and authors only (Trivanovic, Melas, Auderset, Hammer et al.). No abstract was retrievable at harvest. The title indicates an evaluation of existing particle-number periodic-technical-inspection (PN-PTI) instruments - developed and type-approved for diesel vehicles with particle filters - when applied to gasoline exhaust, whose size distribution sits lower and nearer the counters' 23 nm or 10 nm cut. Instruments, test vehicles and error magnitudes are not available; nothing is inferred. Relevance: high - PN-PTI counters are the first mass-produced, low-cost, legally metrological particle counters, and their size-cut behaviour is the same question for community UFP sensing.}
{10.1016/j.jaerosci.2026.106901}

\paperentry{Spinosa volcanic Saharan dust biomass burning aerosol advections Naples 2026}
{Multiple advections of volcanic, Saharan dust and biomass burning aerosols over Naples}
{Atmos Res}
{\textsc{metadata only.} Title and authors only (Spinosa, Manzo, Amoruso, Sannino et al.). No abstract was retrievable at harvest. The title describes overlapping advections of volcanic, Saharan dust and biomass-burning aerosol over Naples; the author group operates lidar there, so a profiling case study is likely but is not confirmed. Instruments, dates and surface PM impacts are not available; nothing is inferred. Relevance: medium - mixed dust/smoke/volcanic layers are the cases where optical low-cost sensors misreport mass most, because refractive index and size change simultaneously.}
{10.1016/j.atmosres.2026.109425}

\par\vskip 0pt plus 10\baselineskip\penalty-200\vskip 0pt plus -10\baselineskip
\band{Occupational \& indoor\hfill 2 records}{Amber}

\paperentry{Martin simplified method portable air cleaner CADR low-cost sensors 2026}
{A simplified method for evaluating fine particulate matter removal performance of portable air cleaners}
{MethodsX}
{Methods paper for community-run CADR testing: release a short-lived nebuliser aerosol in a sealed room, log PM$_{2.5}$ decay with and without the cleaner on low-cost sensors, and compute CADR from the decay-rate difference. Residential low-cost-sensor CADRs ran \textasciitilde{}26\% below laboratory values from research-grade instruments - a conservative bias. Of ten units, a four-filter Corsi-Rosenthal box performed best among DIY builds; a fan shroud added \textasciitilde{}11\%. Across 135 models DIY units cost far less per CADR unit. Limitation: the 26\% gap mixes sensor size-response with room mixing and leakage, and is not decomposed; nebuliser aerosol is not wildfire smoke. Relevance: high - a practical protocol in which the low-cost sensor is the reference, with its bias quantified.}
{10.1016/j.mex.2026.104216}

\paperentry{Hunter e-cigarette aerosol early post-emission particle decay controlled conditions 2026}
{Early post-emission particle decay dynamics of E-cigarette aerosols under controlled conditions}
{Sci Total Environ}
{\textsc{metadata only.} Title and authors only (Hunter, Jacoby, Ferdousi-Rokib, Asa-Awuku et al.). No abstract was retrievable at harvest. The title describes the first seconds-to-minutes decay of e-cigarette aerosol under controlled conditions - evaporation of volatile droplets that shrinks particles rapidly after exhalation or emission. Instruments, conditions and decay constants are not available; nothing is inferred. Relevance: medium - e-cigarette droplets produce large, short-lived spikes on indoor optical PM sensors that mass-based readings misstate because the particles are liquid and evaporating; useful for indoor network event classification.}
{10.1016/j.scitotenv.2026.182379}

\par\vskip 0pt plus 10\baselineskip\penalty-200\vskip 0pt plus -10\baselineskip
\band{Burden, policy \& mitigation\hfill 2 records}{Coral}

\paperentry{Wang climate-adaptive city pilot PM$_{2.5}$ health dividend China 2026}
{Does climate adaptation policy deliver a health dividend? Exposure, modeled health burden, and defensive capacity in China's climate-adaptive city pilots}
{Front Public Health}
{Difference-in-differences on a 2011-2022 panel of 264 prefecture cities, 21 treated by China's 2017 climate-adaptive city pilot. The pilot is associated with PM$_{2.5}$ lower by 3.87 $\mu$g\,m$^{-3}$ (95\% CI -6.49 to -1.25; wild-bootstrap p=0.0067); GEMM translates this into an \textasciitilde{}8.1\% lower modelled burden (-11.32 to -4.76\%) - an accounting result, not observed mortality. Defensive-capacity responses (physicians, beds, insurance) are imprecise. Limitation: the estimate shrinks and loses precision in the matched sample and after pollution-convergence adjustment, so selection of already-improving cities is not excluded; 21 treated units is a small base. Relevance: low-medium - policy evaluation at city scale rests on gridded PM$_{2.5}$ whose own error is not propagated.}
{10.3389/fpubh.2026.1987084}

\paperentry{Wang carbon footprint air-pollutant emissions Beijing 2022 Winter Olympics 2026}
{Linking Carbon Footprints to Air-Pollutant Emissions from Mega Sporting Events: Evidence from the Beijing 2022 Winter Olympics}
{Atmosphere}
{Screening-level inventory converting the verified 713,788 t CO$_2$e footprint of the Beijing 2022 Winter Olympics (2016-2022) into co-emitted pollutants via Chinese pollutant-to-CO$_2$ ratios over eight emission carriers, with Shapley attribution and Monte Carlo/Sobol uncertainty. Construction produced 63\% of CO$_2$e, \textasciitilde{}70\% of 1,335 t NOx and 76\% of 279 t PM$_{2.5}$. Green electricity avoided 31\% of no-measure emissions; offsets reduced no air pollutant. Limitation: emission-side only - ratio-based PM$_{2.5}$ with no dispersion, concentration or exposure step, so there is no health quantity, and national average ratios may misstate construction-site fugitive dust. Relevance: low - construction dust is the source class that site-perimeter low-cost PM monitoring is typically deployed to police.}
{10.3390/atmos17111072}

\clearpage
\section*{\sffamily\bfseries\color{Deep}Corpus register}
\vspace{-1.5mm}{\color{Amber}\rule{\linewidth}{1.1pt}}\vspace{2.5mm}

\input{table_rows}

\vspace{2mm}

\begin{keybox}
{\sffamily\bfseries\footnotesize\color{Deep}Provenance and method}\\[1.2mm]
{\fontsize{7.2}{8.8}\selectfont
Entry window \textbf{30 October 2026}, single day (Friday). Harvester legs: \textbf{PubMed} health
\textbf{16}, sensing \textbf{4}; \textbf{Europe PMC} \textbf{40} (9 AGRICOLA retro-index records dropped);
\textbf{Crossref by ISSN} \textbf{43}; \textbf{arXiv} 150 kB, no in-window entry; \textbf{OpenAlex} HTTP 429. 95 unique,
10 already seen, 85 fresh; PubMed connector (EDAT 30--31 Oct, health and sensing axes separately) 27 PMIDs, 18 bucketed
to 30 Oct, 1 merged, \textbf{6 appended}, \textbf{91 screened}. Abstracts for 22 PubMed records recovered by efetch;
Crossref/Semantic Scholar recovered further abstracts, leaving 19 candidates abstract-free. \textbf{Consensus}
calibration sweep (once for 30--31 October): 10 hits, none new for this day. \textbf{24 admitted}, \textbf{67 rejected}
with logged reasons (5 were EPMC re-surfacings of 29 October rejections; 1 duplicate MethodsX record). \textbf{Trial
watch} 30 October--1 November: 0 registrations updated. \textbf{5 carried on metadata alone}. \textbf{Scholar Gateway}
not called. Abstracts remain the copyright of their publishers.\par}
\end{keybox}

\end{document}
